% Figuras originais para a abertura histórica. Referências: dados/FONTES_CONTEXTO.md.
% As cores distinguem Gate, canal e interconexões; desenhos sem escala.
\newcommand{\linhaHistorica}{%
\begin{tikzpicture}[x=1cm,y=1cm,font=\small]
\draw[seta] (0,0)--(10.3,0);
\foreach \x/\ano/\titulo/\sub in {0.4/1947/Transistor/Bell Labs,2.8/1958--59/Circuito integrado/Kilby e Noyce,5.2/1960--63/MOS e CMOS/Controle por campo,7.6/1971/Microprocessador/Intel 4004,10/Hoje/Sistemas integrados/Múltiplas funções}{%
\fill[blue_theme] (\x,0) circle(.09);
\node[blue_theme,above=3mm,align=center,font=\small\bfseries] at (\x,0){\ano};
\node[below=3mm,align=center,text width=2.25cm,font=\scriptsize] at (\x,0){\textbf{\titulo}\\[1mm]\sub};}
% Dispositivo discreto, circuito, CPU, SoC: motivos conceituais.
\draw[blue_theme,thick] (.05,1.5) rectangle (.75,2.25);
\foreach \x in {.15,.4,.65}{\draw[blue_theme] (\x,1.5)--(\x,1.1);}
\draw[blue_theme,thick] (2,1.2) rectangle (3.6,2.5);
\foreach \x in {2.2,2.7,3.2}{\draw[blue_theme] (\x,1.6)--++(.2,.45)--++(.2,-.45);}
\draw[blue_theme,thick] (4.6,1.4)--(5.8,1.4);\draw[blue_theme,thick] (4.6,1.65)--(5.8,1.65);\draw[blue_theme,line width=4pt] (4.9,2)--(5.5,2);
\draw[blue_theme,thick] (6.8,1.2) rectangle (8.4,2.5);\node at (7.6,1.85){CPU};
\draw[blue_theme,thick] (9.1,1.2) rectangle (10.9,2.5);
\foreach \x/\y in {9.25/1.35,10.05/1.35,9.25/1.95,10.05/1.95}{\fill[blue_theme!20] (\x,\y) rectangle ++(.65,.4);}
\path[use as bounding box] (-.8,-1.3) rectangle (11.2,2.9);
\end{tikzpicture}}

\newcommand{\graficoMoore}{%
\begin{tikzpicture}[x=.86cm,y=.19cm,font=\small]
\draw[->,blue_theme] (0,0)--(9,0) node[right]{anos};
\draw[->,blue_theme] (0,0)--(0,19) node[above,align=center]{quantidade\\relativa};
\foreach \y in {1,2,4,8,16}{\draw[gray!25] (0,\y)--(8,\y);\node[left] at (0,\y){\y};}
\foreach \x in {0,2,4,6,8}{\draw (\x,0)--(\x,-.4) node[below]{\x};}
\draw[blue_theme,very thick,domain=0:8,samples=60] plot (\x,{2^(\x/2)});
\foreach \x/\y in {0/1,2/2,4/4,6/8,8/16}{\fill[blue_theme] (\x,\y) circle[x radius=.07,y radius=.32];}
\node[align=center,fill=white,inner sep=2mm] at (3.6,14){\textbf{Dobrar a cada dois anos}\\$N(t)=N_0\,2^{t/2}$};
\path[use as bounding box] (-1.4,-3) rectangle (10.2,24);
\end{tikzpicture}}

\newcommand{\graficoProcessos}{%
\begin{tikzpicture}[x=.175cm,y=.79cm,font=\scriptsize]
\fill[blue_theme!6] (47,-.1) rectangle (56,4.8);
\draw[->,blue_theme] (0,0)--(57,0);\draw[->,blue_theme] (0,0)--(0,4.8);
\node[anchor=west] at (0,5.15){Referência de processo (nm) · eixo logarítmico};
\foreach \y/\t in {0/1,1/10,2/100,3/{1\,000},4/{10\,000}}{\draw[gray!25] (0,\y)--(56,\y);\node[left] at (0,\y){\t};}
\draw[blue_theme,very thick] (0,4)--(7,3.4771)--(14,3.1761)--(22,2.9031)--(29,2.2553)--(37,1.6532);
\foreach \x/\y/\valor/\ano in {0/4/{10\,000}/1971,14/3.1761/{1\,500}/1985,29/2.2553/180/2000,37/1.6532/45/2008}{%
\fill[blue_theme] (\x,\y) circle[x radius=.45,y radius=.1];
\ifnum\x>0\node[above=1.5mm] at (\x,\y){\valor};\fi\node[below=2mm] at (\x,0){\ano};}
\foreach \x/\y/\nome/\ano in {49/.699/N5/2020,54/.301/N2/2025}{%
\draw[lacuna,thick,fill=white] (\x,\y) circle[x radius=.5,y radius=.11];
\node[above=2mm,lacuna] at (\x,\y){\nome};\node[below=2mm] at (\x,0){\ano};}
\node[align=center,text width=1.7cm] at (51.5,3.4){Nós comerciais\\recentes};
\node[anchor=west] at (0,-1.1){$1\,\mu$m = $1\,000$ nm};
\path[use as bounding box] (-7,-1.5) rectangle (58,5.6);
\end{tikzpicture}}

\newcommand{\litografiaEtapa}[1]{%
\begin{tikzpicture}[x=1cm,y=1cm,font=\small]
\fill[blue_theme!15] (0,0) rectangle (7,.7);\node at (3.5,.32){substrato};
\ifnum#1<3
\fill[blue_theme!65] (0,.7) rectangle (7,.95);
\else
\foreach \x in {0,2.8,5.6}{\fill[blue_theme!65] (\x,.7) rectangle ++(1.4,.25);}
\fi
\ifnum#1=1
\fill[yellow!70] (0,.95) rectangle (7,1.3);
\foreach \x in {1.4,4.2}{\fill[red!12] (\x,.95) rectangle ++(1.4,.35);\fill[red!10] (\x,1.4) rectangle ++(1.4,2.2);\draw[-{Latex},red] (\x+.7,3.5)--(\x+.7,1.45);}
\draw[blue_theme,line width=4pt] (0,3.8)--(1.4,3.8) (2.8,3.8)--(4.2,3.8) (5.6,3.8)--(7,3.8);
\node[left] at (0,3.8){máscara};\node[right,align=left,font=\scriptsize] at (7.2,2.6){Luz projetada\\sobre o resiste};
\else
\foreach \x in {0,2.8,5.6}{\fill[yellow!70] (\x,.95) rectangle ++(1.4,.35);}
\ifnum#1=2
\foreach \x in {2.1,4.9}{\draw[seta] (\x,2.6)--(\x,1.5);}
\node[align=center] at (3.5,3.5){Revelação remove as partes expostas\\de um resiste positivo.};
\else
\foreach \x in {2.1,4.9}{\draw[-{Latex},thick,red] (\x,2.6)--(\x,.85);}
\node[align=center] at (3.5,3.5){Gravação transfere o padrão\\à camada abaixo do resiste.};
\fi\fi
\draw[blue_theme] (0,0) rectangle (7,.7);
\draw[-,blue_theme] (7,.82)--(7.6,.82) node[right,font=\scriptsize]{camada a padronizar};
\draw[-,lacuna] (6.4,1.2)--(7.6,1.7) node[right,font=\scriptsize]{resiste};
\path[use as bounding box] (-1.5,-.4) rectangle (10.1,4.5);
\end{tikzpicture}}

\newcommand{\transistoresModernos}{%
\begin{tikzpicture}[x=1cm,y=1cm,font=\small]
\foreach \x/\nome in {0/Planar,3.8/FinFET,7.6/{GAA · nanosheets}}{%
\node[blue_theme,font=\small\bfseries] at (\x+1.25,3.2){\nome};
\fill[gray!12] (\x,0) rectangle ++(2.5,.5);}
% Cortes transversais ao canal. Dielétrico branco entre Gate e silício.
\fill[ativo!60] (.4,.5) rectangle (2.1,.85);
\fill[blue_theme!70] (.35,1) rectangle (2.15,1.5);
\fill[blue_theme!70] (4.35,.55) rectangle (5.75,2.4);
\fill[white] (4.75,.5) rectangle (5.35,2.05);
\fill[ativo!60] (4.87,.5) rectangle (5.23,1.93);
\fill[blue_theme!70] (8.1,.65) rectangle (9.6,2.55);
\foreach \y in {.9,1.5,2.1}{\fill[white,rounded corners=2pt] (8.3,\y) rectangle ++(1.1,.3);\fill[ativo!60,rounded corners=1pt] (8.4,\y+.08) rectangle ++(.9,.14);}
\node[align=center,font=\scriptsize] at (1.25,-.5){Gate sobre\\a superfície};
\node[align=center,font=\scriptsize] at (5.05,-.5){Gate envolve\\três faces};
\node[align=center,font=\scriptsize] at (8.85,-.5){Gate envolve\\cada canal};
\fill[blue_theme!70] (2,-1.45) rectangle ++(.25,.2);\node[right,font=\scriptsize] at (2.3,-1.35){Gate};
\fill[ativo!60] (5,-1.45) rectangle ++(.25,.2);\node[right,font=\scriptsize] at (5.3,-1.35){canal};
\path[use as bounding box] (-.3,-1.7) rectangle (10.4,3.6);
\end{tikzpicture}}

\newcommand{\pilhaCMOSdois}{%
\begin{tikzpicture}[x=1cm,y=1cm,font=\small]
\foreach \y/\cor/\titulo in {0/gray/Alimentação e acesso externo,1.2/ativo/Lógica densa,2.4/blue_theme/Memória,3.6/lacuna/Interfaces e funções especializadas}{%
\fill[\cor!14] (0,\y)--(5,\y)--(5.8,\y+.45)--(.8,\y+.45)--cycle;
\draw[\cor,thick] (0,\y)--(5,\y)--(5.8,\y+.45)--(.8,\y+.45)--cycle;
\node[anchor=west,font=\scriptsize] at (6,\y+.2){\titulo};}
\foreach \x in {1,2,3,4}{\draw[blue_theme,thick,densely dashed] (\x,.5)--(\x,3.5);}
\node[anchor=west,font=\scriptsize,align=left] at (6,-.9){Conexões 3D densas\\entre camadas funcionais};
\path[use as bounding box] (-.3,-1.4) rectangle (10.6,4.5);
\end{tikzpicture}}
